% Ableitung und Änderungsrate · Die Rate als Funktion · Prüfungs-Fokus Abitur GK – bank/ableitung-und-aenderungsrate/e4.jsonl (zusammenbau v0.6, Prüfungs-Fokus)
\einheitenkopf[e4]{Einheit 4 · Die Rate als Funktion}
\zweigzeile{Abi GK ×5}
\setcounter{aufgabe}{0}

% Anlauf (ohne Prüfkennung)
\begin{aufgabe}{Die Rate als Funktion – Anlauf}
\begin{teile}
\teil $f(t)$ ist die Wassermenge in einem Becken in Litern nach $t$ Minuten. Was beschreibt $f$, und wo ist die größte Zuflussrate zu suchen? \\ \kreuz{Bestand (Höhe, Anzahl, Menge) – größte Rate bei den Nullstellen der zweiten Ableitung} \\ \kreuz{Rate (je Zeiteinheit) – größte Rate bei den Nullstellen der ersten Ableitung}
\teil Die Zuflussrate in ein Becken wird durch $r$ beschrieben ($t$ in Stunden, $r(t)$ in Litern je Stunde); es gilt $r'(t) = (6 - t) \cdot \mathrm{e}^{-0{,}5t}$. Gib die Nullstelle von $r'$ an und nenne den Zeitpunkt der größten Zuflussrate an. t = \leerfeld
\teil Die Temperatur in einem Gewächshaus wird durch $T(t) = -0{,}1t^3 + 1{,}8t^2 + 5$ beschrieben ($t$ in Stunden nach Sonnenaufgang, $T(t)$ in °C, $0 \le t \le 12$). Wann steigt die Temperatur am schnellsten, und wie groß ist die Änderungsrate dann? t = \leerfeld , \leerfeld °C je h
\end{teile}
\end{aufgabe}

\begin{aufgabe}{Die Rate als Funktion – Prüfungsaufgaben}
\begin{teile}
\teil Gegeben ist $f(x) = 0{,}5x^3 - 3x + 2$. Unter den Tangenten an den Graphen von $f$ hat eine die kleinste Steigung. Berechne diese Steigung \hfill (A19). m = \leerfeld
\teil Ein Höhenprofil wird für $0 \le x \le 8$ durch $p(x) = (2x + 1) \cdot \mathrm{e}^{-0{,}4x}$ beschrieben (1 LE = 1 km); $p'(x) = (1{,}6 - 0{,}8x) \cdot \mathrm{e}^{-0{,}4x}$. Weise nach, dass es im Intervall $[2; 8]$ eine Stelle gibt, an der die Steigung des Profils kleiner als $-0{,}33$ ist \hfill (A18).
\teil Die Höhe eines Bambus wird durch $w(t) = -t^3 + 12t^2$ beschrieben ($t$ in Wochen, $w(t)$ in cm, $0 \le t \le 8$); $w'(t) = -3t^2 + 24t$. Berechne die höchste Wachstumsgeschwindigkeit in cm je Woche; die notwendige Bedingung genügt \hfill (A19). \leerfeld[cm] je Woche
\teil Die momentane Förderrate einer Mine wird durch $k'(x) = 0{,}3x \cdot \mathrm{e}^{-0{,}1x}$ beschrieben ($x$ in Jahren seit Beginn der Förderung, $k'(x)$ in Mio. t je Jahr). Berechne den Zeitpunkt der größten Förderrate und ihren Wert \hfill (A26). x = \leerfeld , \leerfeld Mio. t je Jahr
\teil Die Höhe einer Drohne wird für $0 \le t \le 15$ durch $h(t) = 0{,}2t^3 - 6t^2 + 45t$ beschrieben ($t$ in Minuten, $h(t)$ in m); die Phase des Aufstiegs dauert von 0 bis 5 Minuten, $h'(t) = 0{,}6t^2 - 12t + 45$. Bestimme den Zeitraum in dieser Phase, in dem die Steiggeschwindigkeit mindestens $23{,}4$ m je min beträgt \hfill (A22).
\teil Der Füllstand eines Tanks wird für $0 \le t \le 10$ durch $f(t) = -0{,}25t^3 + 3t^2 + 5$ beschrieben ($t$ in Minuten, $f(t)$ in Litern); die Zuflussphase dauert von 0 bis 8 Minuten, $f'(t) = -0{,}75t^2 + 6t$. Bestimme den Zeitraum in dieser Phase, in dem die Zuflussrate mindestens $9$ L je min beträgt \hfill (A22).
\end{teile}
\end{aufgabe}

\begin{aufgabe}{Die Rate als Funktion – Prüfungsaufgaben – weiter}
\begin{teile}
\teil Zwei Bäume wachsen nach $a(t) = 18t^2 - t^3$ und $b(t) = 3t^2$ ($t$ in Jahren, Höhe in cm, $0 \le t \le 12$). Berechne den Zeitpunkt, zu dem der Unterschied der Wachstumsgeschwindigkeiten am größten ist, und diesen Unterschied \hfill (A19).
\teil Zwei Fahrzeuge starten gleichzeitig; ihre Wege sind $s(t) = 0{,}05t^3 + 2t$ und $u(t) = 1{,}5t^2$ ($t$ in Sekunden, Weg in m, $0 \le t \le 15$). Berechne den Zeitpunkt, zu dem der Unterschied ihrer Geschwindigkeiten am größten ist, und diesen Unterschied \hfill (A19).
\teil Die Verkaufszahlen zweier Produkte werden durch $a(t) = t^3 - 9t^2 + 30t$ und $b(t) = 15t$ beschrieben ($t$ in Monaten, verkaufte Stück in Tausend, $0 \le t \le 6$). Bestimme die Zeitpunkte, zu denen der Unterschied der Verkaufsraten am größten ist \hfill (A19).
\teil Die Temperatur eines Getränks im Kühlschrank wird durch $T(x) = 50 \cdot \mathrm{e}^{-0{,}02x} + 6$ beschrieben ($x$ in Minuten, $T(x)$ in °C); $T(x) - 6$ ist die Differenz zur Kühlschranktemperatur, $T'(x) = -\mathrm{e}^{-0{,}02x}$. Aussage: Es gibt eine Zahl $c$, sodass für alle $x$ gilt $T(x) - 6 = c \cdot T'(x)$. Begründe, dass die Aussage im Modell wahr ist \hfill (A26).
\teil Die Menge eines Medikaments im Blut wird durch $m(t) = 80 \cdot \mathrm{e}^{-0{,}25t}$ beschrieben ($t$ in Stunden, $m(t)$ in mg); $m'(t) = -20 \cdot \mathrm{e}^{-0{,}25t}$. Aussage: Zu jedem Zeitpunkt ist die Menge proportional zur momentanen Abbaurate, $m(t) = c \cdot m'(t)$ für alle $t$. Begründe, dass die Aussage im Modell wahr ist \hfill (A26).
\teil Der Ladestand eines Akkus wird durch $L(t) = 100 - 60 \cdot \mathrm{e}^{-0{,}05t}$ beschrieben ($t$ in Minuten, $L(t)$ in Prozent); $L'(t) = 3 \cdot \mathrm{e}^{-0{,}05t}$. Aussage: Es gibt eine Zahl $c$, sodass für alle $t$ gilt $L(t) - 100 = c \cdot L'(t)$. Begründe, dass die Aussage im Modell wahr ist \hfill (A26).
\teil Die Einwohnerzahl einer Stadt wird durch $B(t) = 500 \cdot \mathrm{e}^{0{,}04t} + 200$ beschrieben ($t$ in Jahren, $B(t)$ in Tausend); $B'(t) = 20 \cdot \mathrm{e}^{0{,}04t}$. Aussage: Es gibt eine Zahl $c$, sodass für alle $t$ gilt $B(t) - 200 = c \cdot B'(t)$. Begründe, dass die Aussage im Modell wahr ist, und gib $c$ an \hfill (A26). c = \leerfeld
\end{teile}
\end{aufgabe}

\medskip
%% Merkkasten Einheit 4: 8 Zeilen, Vorgabe höchstens fünf (3.1) – ungekürzt
\uebersichtskasten{Bestand oder Rate? Beschreibt f den Bestand (Höhe, Menge, Anzahl), ist f' die Rate: die größte Rate liegt beim Maximum von f', also bei f'' = 0 mit Vorzeichenwechsel – an der Wendestelle des Bestands. Ist die gegebene Funktion r selbst die Rate (Zufluss je Stunde, Geschwindigkeit), liegt die größte Rate beim Maximum von r, also bei r' = 0. \\ B(t) = 9t² − t³ (Bestand): B'(t) = 18t − 3t², B''(t) = 18 − 6t = 0 ⇔ t = 3; B''' < 0 → größte Rate B'(3) = 27 zum Zeitpunkt 3. \\ Rand: Auf einem Zeitintervall die Kandidaten aus der notwendigen Bedingung und die beiden Randwerte der Rate vergleichen – größter und kleinster Wert können am Rand liegen. \\ B auf [0; 5]: B'(0) = 0, B'(3) = 27, B'(5) = 15 → größte Rate 27 bei 3, kleinste Rate 0 am Rand bei 0. \\ Zeitraum mit Mindestrate: f'(x) = c lösen, beide Lösungen bestimmen, Lösungen außerhalb des betrachteten Bereichs verwerfen; die Länge des Zeitraums ist die Differenz der Lösungen. \\ B'(t) = 15 ⇔ 18t − 3t² = 15 ⇔ t² − 6t + 5 = 0 ⇔ t = 1 oder t = 5 → die Rate ist von 1 bis 5 mindestens 15, der Zeitraum ist 4 lang. \\ Vorzeichen: Eine negative Rate heißt Abnahme des Bestands. Beschreibt die Rate eine Anzahl je Zeit, die nicht negativ werden kann (Eingänge je Stunde), ist das Modell hinter der Nullstelle mit Vorzeichenwechsel ungeeignet. \\ Antwortsatz: Zeitpunkt (mit Uhrzeit, wenn die Zeit so gezählt wird) und Wert der Rate mit Einheit – „Nach 3 Stunden ist die Rate mit 27 Litern je Stunde am größten.“}
